
\documentclass[letterpaper,11pt]{article}
%%%%%%%%%%%%%%%%%%%%%%%%%%%%%%%%%%%%%%%%%%%%%%%%%%%%%%%%%%%%%%%%%%%%%%%%%%%%%%%%%%%%%%%%%%%%%%%%%%%%%%%%%%%%%%%%%%%%%%%%%%%%%%%%%%%%%%%%%%%%%%%%%%%%%%%%%%%%%%%%%%%%%%%%%%%%%%%%%%%%%%%%%%%%%%%%%%%%%%%%%%%%%%%%%%%%%%%%%%%%%%%%%%%%%%%%%%%%%%%%%%%%%%%%%%%%
\usepackage{graphicx}
\usepackage{amsmath,amsfonts,amssymb,amsthm,float}
\usepackage{hyperref}
\usepackage[utf8]{inputenc}
\usepackage[left=2cm, right=2cm, top=2cm, bottom=2cm]{geometry}

\setcounter{MaxMatrixCols}{10}
%TCIDATA{OutputFilter=LATEX.DLL}
%TCIDATA{Version=5.50.0.2953}
%TCIDATA{<META NAME="SaveForMode" CONTENT="1">}
%TCIDATA{BibliographyScheme=BibTeX}
%TCIDATA{LastRevised=Friday, March 27, 2026 00:54:53}
%TCIDATA{<META NAME="GraphicsSave" CONTENT="32">}

\input{tcilatex}
\renewcommand{\baselinestretch}{1.15}
\setlength{\parindent}{0pt}
\setlength{\parskip}{0.5\baselineskip}
\pretolerance=2000 \tolerance=3000
\renewcommand{\abstractname}{Resumen}

\begin{document}

\title{Pr\'{a}ctica 2:\ Caos en sistemas biol\'{o}gicos}
\author{Arellano Salones Bricia Idalia $22211746$ \\
%EndAName
Departamento de Ingenier\'{\i}a El\'{e}ctrica y Electr\'{o}nica\\
Tecnol\'{o}gico Nacional de M\'{e}xico / Instituto Tecnol\'{o}gico de Tijuana%
}
\maketitle

\noindent \textbf{Palabras clave: }C\'{a}ncer e inmunoterapia; Caos;
Estabilidad; Simulaciones num\'{e}ricas; Sistema no lineal.

\noindent

\noindent Correo: \textbf{L22211746@tectijuana.edu.mx}

\noindent \noindent Carrera: \textbf{Ingenier\'{\i}a Biom\'{e}dica }

\noindent Asignatura: \textbf{Biolog\'{\i}a de Sistemas}

\noindent Profesor: \href{https://biomath.xyz/}{\textbf{Dr. Paul Antonio
Valle Trujillo}} (paul.valle@tectijuana.edu.mx)

\section{Modelo matem\'{a}tico}

El modelo matem\'{a}tico que describe la din\'{a}mica de tres poblaciones
celulares est\'{a} descrito por las siguientes tres ecuaciones diferenciales
ordinarias de primer orden:%
\begin{eqnarray}
\dot{x} &=&r_{1}x\left( 1-b_{1}x\right) -a_{12}xy-a_{13}xz,  \label{x} \\
\dot{y} &=&r_{2}y\left( 1-b_{2}y\right) -a_{21}xy,  \label{y} \\
\dot{z} &=&\left( r_{3}-a_{31}\right) xz-d_{3}z+\rho _{i},  \label{z}
\end{eqnarray}%
donde todos los par\'{a}metros son positivos y el par\'{a}metro de control $%
\rho _{i}\geq 0.$

\section{Positividad del sistema}

Con base en el trabajo de De Leenheer and Aeyels, se realiza el siguiente c%
\'{a}lculo para establecer la positividad del sistema: 
\begin{eqnarray*}
\left. \dot{x}\right\vert _{x=0} &=&r_{1}\left( 0\right) \left(
1-b_{1}\left( 0\right) \right) -a_{12}\left( 0\right) y-a_{13}\left(
0\right) z=0, \\
\left. \dot{y}\right\vert _{y=0} &=&r_{2}\left( 0\right) \left(
1-b_{2}\left( 0\right) \right) -a_{21}x\left( 0\right) =0, \\
\left. \dot{z}\right\vert _{z=0} &=&\left( r_{3}-a_{31}\right) x\left(
0\right) -d_{3}\left( 0\right) +\rho _{i}=\rho _{i},
\end{eqnarray*}%
por lo tanto, se establece el siguiente resultado:

\noindent \textbf{Resultado I. Positividad de las soluciones. }Las
soluciones $\left( x\left( t\right) ,y\left( t\right) ,z\left( t\right)
\right) $ y \ semitrayectorias positivas $\left( \Gamma ^{+}\right) $ del
sistema (\textit{\ref{x}})-(\textit{\ref{z}}) ser\'{a}n positivamente
invariantes y oara cada condici\'{o}n inicial no negativa $x\left( 0\right)
,y\left( 0\right) ,z\left( 0\right) \geq 0$, se localizar\'{a}n en el
siguiente dominio:\ 
\begin{equation*}
\mathbf{R}_{+,0}^{3}=\left\{ x\left( t\right) ,y\left( t\right) ,z\left(
t\right) \geq 0\right\} .
\end{equation*}

\section{Puntos de equilibrio libre de c\'{e}lulas patol\'{o}gicas}

Con el objetivo de determinar condiciones para la eliminaci\'{o}n de la
poblaci\'{o}n de c\'{e}lulas patal\'{o}gicas, se calculan los equilibrio
donde $x^{\ast }=0,$ es decir,%
\begin{eqnarray*}
x &=&0, \\
r_{2}y\left( 1-b_{2}y\right) -a_{21}xy &=&0, \\
\left( r_{3}-a_{31}\right) xz-d_{3}z+\rho _{i} &=&0,
\end{eqnarray*}%
simplificando se obtiene lo siguiente\ 

$\func{assume}\left( r_{2},\func{positive}\right) =\allowbreak \left(
0,\infty \right) $

$\func{assume}\left( b_{2},\func{positive}\right) =\allowbreak \left(
0,\infty \right) $

$\func{assume}\left( d_{3},\func{positive}\right) =\allowbreak \left(
0,\infty \right) $%
\begin{eqnarray*}
r_{2}y\left( 1-b_{2}y\right) &=&0 \\
-d_{3}z+\rho _{i} &=&0
\end{eqnarray*}%
entonces los equilibrios libres de c\'{e}lulas patol\'{o}gicas son los
siguientes:%
\begin{eqnarray*}
\left( x_{6}^{\ast },y_{6}^{\ast },z_{6}^{\ast }\right) &=&\left( 0,0,\frac{%
\rho _{i}}{d_{3}}\right) , \\
\left( x_{4}^{\ast },y_{4}^{\ast },z_{4}^{\ast }\right) &=&\left( 0,\frac{1}{%
b_{2}},\frac{\rho _{i}}{d_{3}}\right) .
\end{eqnarray*}

\section{Estabilidad local}

Ahora, con el prop\'{o}sito de establecer condiciones para asegurar la
estabilidad asint\'{o}tica local del equilibrio $\left( x_{4}^{\ast
},y_{4}^{\ast },z_{4}^{\ast }\right) $, se calcula la matriz Jacobiana:%
\begin{equation*}
J=\left[ 
\begin{array}{ccc}
-a_{12}\,y-a_{13}\,z-r_{1}\,{\left( b_{1}\,x-1\right) }-b_{1}\,r_{1}\,x & 
-a_{12}\,x & -a_{13}\,x \\ 
-a_{21}\,y & -a_{21}\,x-r_{2}\,{\left( b_{2}\,y-1\right) }-b_{2}\,r_{2}\,y & 
0 \\ 
-z\,{\left( a_{31}-r_{3}\right) } & 0 & -d_{3}-x\,{\left(
a_{31}-r_{3}\right) }%
\end{array}%
\right] ,
\end{equation*}%
y se eval\'{u}a como se indica a continuaci\'{o}n:

\begin{eqnarray*}
\left. J\right\vert _{\left( x_{4}^{\ast },y_{4}^{\ast },z_{4}^{\ast
}\right) } &=&\left[ 
\begin{array}{ccc}
-a_{12}\,y-a_{13}\,z-r_{1}\,{\left( b_{1}\,x-1\right) }-b_{1}\,r_{1}\,x & 
-a_{12}\,x & -a_{13}\,x \\ 
-a_{21}\,y & -a_{21}\,x-r_{2}\,{\left( b_{2}\,y-1\right) }-b_{2}\,r_{2}\,y & 
0 \\ 
-z\,{\left( a_{31}-r_{3}\right) } & 0 & -d_{3}-x\,{\left(
a_{31}-r_{3}\right) }%
\end{array}%
\right] _{x=0,y=\frac{1}{b_{2}},z=\frac{\rho _{i}}{d_{3}}}, \\
\left. J\right\vert _{\left( x_{4}^{\ast },y_{4}^{\ast },z_{4}^{\ast
}\right) } &=&\left[ 
\begin{array}{ccc}
r_{1}-\frac{a_{12}}{b_{2}}-\frac{a_{13}\rho _{i}}{d_{3}} & 0 & 0 \\ 
-\frac{a_{21}}{b_{2}} & -r_{2} & 0 \\ 
\frac{\rho _{i}}{d_{3}}\left( r_{3}-a_{31}\right) & 0 & -d_{3}%
\end{array}%
\right] ,
\end{eqnarray*}
:

con base en la estructura de la matriz (matriz triangular),\ los valores
propios est\'{a}n dados por: 
\begin{eqnarray*}
\lambda _{1} &=&r_{1}-\frac{a_{12}}{b_{2}}-\frac{a_{13}\rho _{i}}{d_{3}}, \\
\lambda _{2} &=&-r_{2}, \\
\lambda _{3} &=&-d_{3},
\end{eqnarray*}%
por lo tanto, para asegurar que el equilibrio es $\left( x_{4}^{\ast
},y_{4}^{\ast },z_{4}^{\ast }\right) $ es local asint\'{o}ticamente estable,
se establece la siguiente condici\'{o}n:\ 
\begin{equation*}
r_{1}-\frac{a_{12}}{b_{2}}-\frac{a_{13}\rho _{i}}{d_{3}}<0,
\end{equation*}%
lo cual se resuelve para el par\'{a}metro de tratamiento $\rho _{i}$ 
\begin{equation*}
\rho _{i}>\left( r_{1}-\frac{a_{12}}{b_{2}}\right) \frac{d_{3}}{a_{13}}.
\end{equation*}

\end{document}
